%=========================================================
% C++ Chapter 9 Lab Handout (Verbatim Korean-safe)
% Topic: Virtual functions & Abstract classes (가상 함수와 추상 클래스)
% Compile with XeLaTeX
%=========================================================
\documentclass[11pt,a4paper]{article}

\usepackage[margin=1in]{geometry}
\usepackage{fontspec}
\usepackage{kotex}
\usepackage{hyperref}
\usepackage{enumitem}
\usepackage{fancyhdr}
\usepackage{titlesec}
\usepackage{xcolor}
\usepackage{tcolorbox}
\usepackage{fvextra}

% ---------- Fonts ----------
% If D2Coding isn't available, try: NanumGothicCoding, Noto Sans Mono CJK KR
\setmonofont{D2Coding}[Scale=MatchLowercase]

% ---------- Colors ----------
\definecolor{CodeBg}{RGB}{248,248,248}
\definecolor{CodeFrame}{RGB}{220,220,220}
\definecolor{Accent}{RGB}{40,80,160}

% ---------- Header ----------
\pagestyle{fancy}
\fancyhf{}
\lhead{\textbf{C++ 9장 실습}}
\rhead{\textbf{가상 함수와 추상 클래스}}
\cfoot{\thepage}

\titleformat{\section}{\Large\bfseries\color{Accent}}{\thesection}{0.6em}{}
\titleformat{\subsection}{\large\bfseries}{\thesubsection}{0.6em}{}

% ---------- Boxes ----------
\newtcolorbox{GoalBox}{
  colback=CodeBg,
  colframe=CodeFrame,
  arc=3pt,
  boxrule=0.8pt,
  left=10pt,right=10pt,top=8pt,bottom=8pt
}
\newcommand{\filetag}[1]{\texttt{\color{Accent}#1}}

% ---------- Code block environment (Unicode-safe) ----------
\DefineVerbatimEnvironment{cppcode}{Verbatim}{
  fontsize=\small,
  frame=single,
  framerule=0.6pt,
  rulecolor=\color{CodeFrame},
  numbers=left,
  numbersep=8pt,
  baselinestretch=1,
  breaklines=true,
  breaksymbolleft=,
  breaksymbolright=,
  commandchars=\\\{\},
  xleftmargin=1.5em
}

\begin{document}

\begin{center}
  {\LARGE \textbf{C++ 9장 실습: 가상 함수와 추상 클래스}}\\
  \vspace{6pt}
  {\large (overriding / virtual / dynamic binding / virtual destructor / pure virtual / abstract class)}
\end{center}

\vspace{8pt}

\section*{학습 목표(슬라이드 요약)}
\begin{GoalBox}
\begin{enumerate}[leftmargin=18pt]
  \item 상속에서 함수 재정의(redefine)를 이해한다.
  \item 가상 함수와 오버라이딩, 동적 바인딩의 개념을 이해한다.
  \item 가상 소멸자의 중요성을 이해한다.
  \item 가상 함수를 활용하여 프로그램을 작성할 수 있다.
  \item 순수 가상 함수와 추상 클래스를 이해하고 작성할 수 있다.
\end{enumerate}
\end{GoalBox}

\section*{실습 운영 규칙}
\begin{itemize}[leftmargin=18pt]
  \item 모든 코드는 \texttt{main()} 포함, \textbf{독립 실행} 가능합니다.
  \item (에러 확인) 표시는 주석을 풀면 일부러 컴파일 에러가 나도록 만든 학습용 코드입니다.
  \item \textbf{9장의 핵심은 “Base* 로 호출했는데 Derived 함수가 실행되는가?”} 입니다.
  \item 삭제(delete)가 들어가는 예제는 \textbf{가상 소멸자}를 반드시 확인하세요.
\end{itemize}

\tableofcontents
\newpage

%=========================================================
\section{가상 함수가 없을 때: 단순 함수 재정의(redefine) (예제 9-1)}
\begin{GoalBox}
\textbf{핵심}
\begin{itemize}[leftmargin=18pt]
  \item 파생 클래스에서 같은 이름의 함수를 만들면 ``함수 재정의''가 됩니다.
  \item 하지만 Base 포인터로 호출하면, \textbf{정적 바인딩}으로 Base 함수가 호출됩니다.
\end{itemize}
\end{GoalBox}

\noindent\filetag{L1\_RedefineNoVirtual.cpp}
\begin{cppcode}
#include <iostream>
using namespace std;

class Base {
public:
    void f() { cout << "Base::f() called\n"; }
};

class Derived : public Base {
public:
    void f() { cout << "Derived::f() called\n"; }
};

int main() {
    Derived d;
    Derived* pDer = &d;
    pDer->f();             // Derived::f()

    Base* pBase = pDer;    // 업캐스팅
    pBase->f();            // Base::f() (가상 함수가 아니면 정적 바인딩)
}
\end{cppcode}

%=========================================================
\section{가상 함수: 오버라이딩 + 동적 바인딩 (예제 9-2)}
\begin{GoalBox}
\textbf{핵심}
\begin{itemize}[leftmargin=18pt]
  \item \texttt{virtual}은 ``동적 바인딩 지시어''입니다.
  \item Base 포인터로 호출해도, 실제 객체가 Derived면 \textbf{Derived 함수가 실행}됩니다.
\end{itemize}
\end{GoalBox}

\noindent\filetag{L2\_VirtualOverriding.cpp}
\begin{cppcode}
#include <iostream>
using namespace std;

class Base {
public:
    virtual void f() { cout << "Base::f() called\n"; }
};

class Derived : public Base {
public:
    void f() override { cout << "Derived::f() called\n"; } // override는 실수 방지용(권장)
};

int main() {
    Derived d;
    Base* pBase = &d;  // 업캐스팅
    pBase->f();        // 동적 바인딩 -> Derived::f()
}
\end{cppcode}

%=========================================================
\section{동적 바인딩 예제: Base 멤버 함수가 내부에서 virtual 호출}
\begin{GoalBox}
\textbf{핵심}
\begin{itemize}[leftmargin=18pt]
  \item Base의 \texttt{paint()}는 Base 함수이지만, 내부에서 \texttt{draw()}(virtual)를 호출합니다.
  \item 그래서 \texttt{Shape* p = new Circle(); p->paint();}에서 \textbf{Circle::draw()}가 실행됩니다.
\end{itemize}
\end{GoalBox}

\noindent\filetag{L3\_PaintDrawDynamic.cpp}
\begin{cppcode}
#include <iostream>
using namespace std;

class Shape {
public:
    void paint() { draw(); } // Base 멤버 함수
    virtual void draw() { cout << "Shape::draw() called\n"; }
};

class Circle : public Shape {
public:
    void draw() override { cout << "Circle::draw() called\n"; }
};

int main() {
    Shape* p1 = new Shape();
    p1->paint();  // Shape::draw()
    delete p1;

    Shape* p2 = new Circle(); // 업캐스팅
    p2->paint();  // Circle::draw() (동적 바인딩)
    delete p2;
}
\end{cppcode}

%=========================================================
\section{상속이 반복될 때: Base/Derived/GrandDerived (예제 9-3)}
\begin{GoalBox}
\textbf{핵심}
\begin{itemize}[leftmargin=18pt]
  \item virtual 함수는 ``가장 아래(가장 구체) 클래스''의 오버라이딩이 선택됩니다.
  \item Base*, Derived*, GrandDerived* 모두 같은 객체를 가리키면, 호출 결과는 동일할 수 있습니다.
\end{itemize}
\end{GoalBox}

\noindent\filetag{L4\_GrandDerivedDynamic.cpp}
\begin{cppcode}
#include <iostream>
using namespace std;

class Base {
public:
    virtual void f() { cout << "Base::f() called\n"; }
};
class Derived : public Base {
public:
    void f() override { cout << "Derived::f() called\n"; }
};
class GrandDerived : public Derived {
public:
    void f() override { cout << "GrandDerived::f() called\n"; }
};

int main() {
    GrandDerived g;
    Base* bp = &g;
    Derived* dp = &g;
    GrandDerived* gp = &g;

    bp->f();  // GrandDerived::f()
    dp->f();  // GrandDerived::f()
    gp->f();  // GrandDerived::f()
}
\end{cppcode}

%=========================================================
\section{범위 지정 연산자(::)로 Base의 virtual 호출하기 (예제 9-4)}
\begin{GoalBox}
\textbf{핵심}
\begin{itemize}[leftmargin=18pt]
  \item \texttt{p->draw()}는 동적 바인딩으로 파생 draw()가 호출될 수 있습니다.
  \item \texttt{p->Shape::draw()}는 \textbf{정적 바인딩}으로 Base의 draw()를 강제로 호출합니다.
\end{itemize}
\end{GoalBox}

\noindent\filetag{L5\_ScopeResolutionVirtual.cpp}
\begin{cppcode}
#include <iostream>
using namespace std;

class Shape {
public:
    virtual void draw() { cout << "--Shape--"; }
};

class Circle : public Shape {
public:
    void draw() override {
        Shape::draw();     // (선택) Base 기능을 먼저 실행하고
        cout << "Circle\n"; // 기능 추가
    }
};

int main() {
    Circle circle;
    Shape* pShape = &circle;
    pShape->draw();           // 동적 바인딩 -> Circle::draw()
    pShape->Shape::draw();    // 정적 바인딩 -> Shape::draw()
    cout << "\n";
}
\end{cppcode}

%=========================================================
\section{가상 소멸자: Base*로 delete 할 때 필수 (예제 9-5)}
\begin{GoalBox}
\textbf{핵심}
\begin{itemize}[leftmargin=18pt]
  \item \textbf{Base 포인터로 Derived 객체를 delete}할 가능성이 있으면, Base 소멸자는 반드시 virtual이어야 합니다.
  \item 그렇지 않으면 Derived 소멸자가 호출되지 않아 자원 해제가 누락될 수 있습니다.
\end{itemize}
\end{GoalBox}

\noindent\filetag{L6\_VirtualDestructor.cpp}
\begin{cppcode}
#include <iostream>
using namespace std;

class Base {
public:
    virtual ~Base() { cout << "~Base()\n"; }
};

class Derived : public Base {
public:
    ~Derived() override { cout << "~Derived()\n"; }
};

int main() {
    Derived* dp = new Derived();
    Base* bp = new Derived();

    delete dp; // ~Derived -> ~Base
    delete bp; // (중요) Base*로 delete 해도 ~Derived -> ~Base
}
\end{cppcode}

%=========================================================
\section{순수 가상 함수와 추상 클래스}
\begin{GoalBox}
\textbf{핵심}
\begin{itemize}[leftmargin=18pt]
  \item \texttt{virtual void draw() = 0;} 처럼 선언하면 \textbf{순수 가상 함수}.
  \item 순수 가상 함수를 1개라도 가진 클래스는 \textbf{추상 클래스}가 되어 객체 생성 불가.
  \item 대신 \textbf{포인터/참조}는 만들 수 있고, 파생 클래스에서 구현하면 사용 가능.
\end{itemize}
\end{GoalBox}

\noindent\filetag{L7\_PureVirtualIntro.cpp}
\begin{cppcode}
#include <iostream>
using namespace std;

class Shape {
public:
    virtual void draw() = 0; // 순수 가상 함수
    void paint() { draw(); } // 순수 가상 함수도 호출은 가능(실행은 파생이 담당)
};

class Circle : public Shape {
public:
    void draw() override { cout << "Circle\n"; }
};

int main() {
    // Shape s;             // (에러 확인) 추상 클래스 객체 생성 불가
    Shape* p = new Circle();  // OK (포인터는 가능)
    p->paint();               // Circle::draw()
    delete p;
}
\end{cppcode}

%=========================================================
\section{실습 9-6: 추상 클래스 Calculator 구현 (슬라이드 예제 9-6)}
\begin{GoalBox}
\textbf{문제}
\begin{itemize}[leftmargin=18pt]
  \item 추상 클래스 \texttt{Calculator}를 상속받아 \texttt{GoodCalc}를 구현한다.
  \item 3개의 순수 가상 함수를 모두 구현해야 객체 생성 가능.
\end{itemize}
\end{GoalBox}

\noindent\filetag{L8\_GoodCalc.cpp}
\begin{cppcode}
#include <iostream>
using namespace std;

class Calculator {
public:
    virtual int add(int a, int b) = 0;
    virtual int subtract(int a, int b) = 0;
    virtual double average(int a[], int size) = 0;
    virtual ~Calculator() {} // (권장) 추상 클래스도 가상 소멸자
};

class GoodCalc : public Calculator {
public:
    int add(int a, int b) override { return a + b; }
    int subtract(int a, int b) override { return a - b; }
    double average(int a[], int size) override {
        double sum = 0;
        for (int i = 0; i < size; i++) sum += a[i];
        return sum / size;
    }
};

int main() {
    int a[] = {1,2,3,4,5};
    Calculator* p = new GoodCalc();
    cout << p->add(2, 3) << endl;
    cout << p->subtract(2, 3) << endl;
    cout << p->average(a, 5) << endl;
    delete p;
}
\end{cppcode}

%=========================================================
\section{실습 9-7: 템플릿 메서드 패턴 느낌의 추상 클래스 (슬라이드 예제 9-7)}
\begin{GoalBox}
\textbf{핵심}
\begin{itemize}[leftmargin=18pt]
  \item 기본 클래스는 \texttt{run()}의 전체 흐름을 제공하고,
        파생 클래스는 \texttt{calc()}만 구현한다.
  \item 이런 구조는 ``공통 흐름 + 부분 구현''을 깔끔하게 분리한다.
\end{itemize}
\end{GoalBox}

\noindent\filetag{L9\_AdderSubtractor.cpp}
\begin{cppcode}
#include <iostream>
using namespace std;

class Calculator {
    void input() {
        cout << "정수 2 개를 입력하세요>> ";
        cin >> a >> b;
    }
protected:
    int a, b;
    virtual int calc(int a, int b) = 0; // 순수 가상 함수
public:
    void run() {
        input();
        cout << "계산된 값은 " << calc(a, b) << endl;
    }
    virtual ~Calculator() {}
};

class Adder : public Calculator {
protected:
    int calc(int a, int b) override { return a + b; }
};

class Subtractor : public Calculator {
protected:
    int calc(int a, int b) override { return a - b; }
};

int main() {
    Adder adder;
    Subtractor subtractor;
    adder.run();
    subtractor.run();
}
\end{cppcode}

%=========================================================
\section*{마무리 체크리스트}
\begin{GoalBox}
\begin{itemize}[leftmargin=18pt]
  \item virtual이 없을 때 Base*로 호출하면 왜 Base 함수가 실행되는가(정적 바인딩)?
  \item virtual이 있으면 Base*로 호출해도 왜 Derived 함수가 실행되는가(동적 바인딩)?
  \item 오버라이딩 성공 조건(이름/매개변수/리턴 타입 일치)을 설명할 수 있는가?
  \item \texttt{override}를 쓰면 어떤 실수를 줄일 수 있는가?
  \item Base의 가상 함수를 \texttt{Base::func()}로 강제 호출하면 무슨 효과(정적 바인딩)가 있는가?
  \item Base*로 delete할 가능성이 있으면 왜 Base 소멸자를 virtual로 해야 하는가?
  \item 순수 가상 함수와 추상 클래스의 의미(객체 생성 불가, 포인터 가능)를 말할 수 있는가?
\end{itemize}
\end{GoalBox}

\end{document}
